\documentclass[output=paper]{langscibook}
\ChapterDOI{10.5281/zenodo.22880256}
\title{Introduction}
\author{Ulrike Mosel}
\abstract{}
\begin{document}
\maketitle\label{chap:15}
The first documents of Teop architecture seem to be a photograph of a village at the Teop Harbour that was taken by Richard Parkinson, see \figref{fig:15_1}. and the photographs by Schaffrath around 1900\footnote{\cite{Parkinson:1907kx}: Fig. 75, ``Dorfszene am Ernst-Günther-Hafen (Nord Bougainville)''; \cite{Schultze:1909kw}: Fig. 38, ``Typisches Stranddorf auf Bougainville''; Fig. 51, ``Typischer Wohnhausbau auf Bougainville.'' K.J. Schaffrath worked as a construction engineer on Bougainville Island for the colonial German government.}, see \figref{fig:15_2} and \figref{fig:15_3}.
\begin{figure}[H]
    \captionsetup[subfigure]{oneside,margin={0cm,0cm}}
    \subfloat[]{\includegraphics[trim=0 0 0 70,clip,width=.35\linewidth]{figures/15_1.jpg}\label{fig:15_1}}
    \subfloat[]{\includegraphics[trim=0 0 0 0,clip,width=.36\linewidth]{figures/15_2.jpg}\label{fig:15_2}}
    \subfloat[]{\includegraphics[trim=0 0 0 0,clip,width=.29\linewidth]{figures/15_3.jpg}\label{fig:15_3}}
    \caption{(a) A village at Teop Harbour, (b) `Typical beach village on Bougainville', (c) `Typical construction of a dwelling house on Bougainville'}
\end{figure}
The anthropologist Beatrice Blackwood, who did fieldwork in Buka and Northern Bougainville for 12 months in 1929-1930, reports that the houses of the first type (\figref{fig:15_1}) were predominant ``in most places of Bougainville'' where they were ``built on bare earth, thatched with sago-palm leaves in the shape of a broad gothic arch, with the roof almost touching the ground on both sides'' (\cite[22]{Blackwood:1935im}). The construction of houses raised on piles was introduced by missionaries (\figref{fig:15_3}). When I was doing my last fieldwork in 2014, houses of the second and the third type were still in use, see \figref{fig:15_4} and \figref{fig:15_5}.
\begin{figure}[H]
    \captionsetup[subfigure]{oneside,margin={0cm,0cm}}
    \captionsetup{justification=centering}
    \subfloat[]{\includegraphics[trim=0 0 0 0,clip,width=.53\linewidth]{figures/15_4.jpg}\label{fig:15_4}}
    \subfloat[]{\includegraphics[trim=0 18 0 0,clip,width=.47\linewidth]{figures/15_5.jpg}\label{fig:15_5}}
    \caption{(a) \textit{A inu vaa putaa} -- A house on the ground,\\ (b) \textit{A inu vaa gina} -- lit. A house on stilts}
\end{figure}
The first type of house (\figref{fig:15_1}) is remembered by the Teop people as \textit{tagoana} `the chief’s house', see \figref{fig:15_6}, and is used as a ceremonial hut during large feasts, see \figref{fig:15_7}. \textcite[383-388]{Blackwood:1935im} describes this traditional house type with two excellent photographs and an index of terms in the Kurtatchi language that denote plant materials and parts of the house. Some of them are etymologically related to the corresponding Teop terms, e.g. \textit{inu} `hut'. The Teop Language Corpus only contains a few mentions of \textit{tagoana} in the sense of `house of a chief' and `ceremonial hut'.
\begin{figure}[H]
    \captionsetup[subfigure]{oneside,margin={0cm,0cm}}
    \captionsetup{justification=centering}
    \subfloat[]{\includegraphics[trim=0 0 0 0,clip,width=.39\linewidth]{figures/15_6.jpg}\label{fig:15_6}}
    \subfloat[]{\includegraphics[trim=0 0 0 0,clip,width=.39\linewidth]{figures/15_7.jpg}\label{fig:15_7}}
    \subfloat[]{\includegraphics[trim=0 0 0 10,clip,width=.22\linewidth]{figures/15_8.jpg}\label{fig:15_8}}
    \caption{(a) \textit{Tagoana} -- House of a chief, (b) Ceremonial hut,\\ (c) Inside the ceremonial hut}
\end{figure}
\textcite[294-295]{Shoffner:1976eg} whose description of the traditional type of house is similar to that of Blackwood mentions that ``this type of house is built very occasionally today for ceremonial purposes but never to live in''.

In the following chapters Enoch Horai Magum describes how in 2005 \textit{inu} \textit{gorogorohu} `sleeping houses' were exclusively build from local plant materials, but at that time houses were also built of a combination of local and imported materials. \textcite[295]{Shoffner:1976eg} reports, ``The roof is most commonly made of sago thatching although today, more people are acquiring metal roofs partially because they are necessary for the collection of water in tanks or any type of receptacle available.''

\figref{fig:15_9} and \figref{fig:15_10} show a house with a roof made of corrugated iron, modern windows, stairs and a door, but its walls are made of woven bamboos in a traditional way, see Chapter \ref{chap:19}.
\begin{figure}[H]
    \captionsetup[subfigure]{oneside,margin={0cm,0cm}}
    \subfloat[]{\includegraphics[trim=0 0 0 17,clip,width=.58\linewidth]{figures/15_9.jpg}\label{fig:15_9}}
    \subfloat[]{\includegraphics[trim=0 0 0 0,clip,width=.42\linewidth]{figures/15_10.jpg}\label{fig:15_10}}
    \caption{\textit{A inu vaa nomana} -- A modern house}
\end{figure}
In contrast to the European dwelling house, the Teop `sleeping house' does not have a kitchen, toilet or bathroom. 
\end{document}